The problem of constructing objective or noninformative priors has a long history in Bayesian statistics.
Among the most influential approaches are Jeffreys priors, reference priors, and maximum-entropy based
methods~\cite{Jeffreys1946, Bernardo1979, jaynes1957}.

Jeffreys introduced a general rule for constructing invariant noninformative priors based on the Fisher information
matrix~\cite{Jeffreys1946}.
The resulting priors are widely used because of their parameterization invariance and their connection to asymptotic
likelihood theory.
However, Jeffreys priors are derived from local geometric properties of the parameter space and do not directly
address the problem of representing a set of admissible probability distributions by a single distribution.

The closest connection to the present work is Bernardo's reference prior
methodology~\cite{BERNARDO200517, Bernardo1979}.
Reference priors are constructed by maximizing the expected information gained from the data and are intended to
represent a state of minimal prior information.
In Section~\ref{subsec:parametric}, we show that central priors satisfy a fixed-point condition closely related to
the reference prior construction.
In particular, for a single parameter and a single set of distributions, the resulting central prior coincides with
the corresponding reference prior.

Several authors have also studied objective Bayesian inference through information-theoretic
principles~\cite{jaynes1957, 10.1214/aoms/1177729694}.
Examples include maximum entropy methods, minimum discrimination information principles, and related
divergence-based approaches.
These methods typically select a distribution by optimizing an information criterion subject to constraints.
In contrast, the central distribution is defined as a minimax representative of a set or hierarchy of distributions,
minimizing the worst-case Kullback–Leibler divergence to the admissible models.

The present work is also related to robust Bayesian statistics and imprecise
probability~\cite{riosinsua2000robust, walley1991statistical}.
In these frameworks uncertainty is represented by sets of probability distributions rather than a single model.
Existing approaches generally focus on performing inference directly with sets of distributions.
The central distribution framework instead seeks a single representative distribution that summarizes the available
set-valued or hierarchical uncertainty while preserving its information-theoretic structure.

The central distribution framework differs from these approaches in two respects.
First, it is defined as a minimax-KL representative of a set, or more generally a hierarchy, of probability
distributions.
Second, it provides a mechanism for converting set-valued and hierarchical uncertainty into a single probability
distribution suitable for Bayesian inference.
While related to objective Bayesian inference, robust Bayesian analysis, and information-theoretic methods, its
primary goal is not the direct specification of a prior distribution for parameters or inference with sets of
distributions, but rather the construction of a representative distribution that preserves the informational
structure of the underlying uncertainty.
